\chapter{Risk in Production: Greeks, P\&L Explain, and xVA}
\label{ch:risk_production}
\chaptermark{Risk in production}

This chapter takes the single-option Greeks of
\cref{ch:greeks_hedging} into production: \emph{portfolio-scale}
Greeks across thousands of positions, \emph{P\&L Explain} as the
Strat's morning ritual, and the \emph{xVA alphabet soup} (CVA, DVA,
FVA, KVA) accounting for what textbook Black--Scholes ignores:
counterparty default, own default, funding costs, regulatory
capital.  The new content is \emph{scale} and \emph{accountability}:
a 10k-trade desk cannot reprice by hand; a Strat who cannot explain
yesterday's P\&L to four significant figures is unemployed by
Friday; a pricing model valuing a derivative at \$1{,}000{,}000
while ignoring \$30k CVA, \$10k FVA, and \$8k KVA is wrong by 5\%
on a product whose bid--offer is 2\%.

\begin{backgroundbox}[Prerequisites]
From earlier chapters:
\begin{itemize}
  \item \Cref{ch:greeks_hedging}: Black--Scholes Greeks for a single
    European option (delta, gamma, vega, theta, rho), and the
    first-order hedging identity $\Delta V \approx \Delta \cdot
    \Delta S$.
  \item \Cref{ch:kelly_risk}: Sharpe ratio, drawdown, VaR, and the
    idea that risk is a \emph{distribution} of P\&L outcomes, not a
    point estimate.
  \item \Cref{ch:fixed_income_foundations}: zero curves and bucketed
    rate risk: the object we called a \emph{key-rate duration}
    generalises to the KRD vector used here.
  \item \Cref{ch:credit_derivatives}: hazard rates $h(t)$, survival
    $S(t) = e^{-\int_0^t h(u)du}$, and the expected-loss formula
    $\mathrm{EL} = (1-R)\,\mathrm{PD}$.  CVA is the
    counterparty-exposure analogue.
\end{itemize}
\end{backgroundbox}


%% ================================================================
\section{Portfolio-scale Greeks}
\label{sec:portfolio_greeks}
%% ================================================================

A real options book holds tens of thousands of positions across
hundreds of underlyings, dozens of expiries, and a range of strikes.
The risk system's job is to \emph{aggregate} the Greeks into a
handful of numbers a trader can act on.

\subsection{Linear aggregation of Greeks}

Greeks of a portfolio are additive in the \emph{dollar} convention,
not in the \emph{per-share} convention.  Let the book hold $n_i$
contracts of option $i$ on underlying $U(i)$, with per-contract
delta $\delta_i$, gamma $\gamma_i$, vega $\nu_i$, and let $m$ be the
contract multiplier (100 shares for US equity options).  The
\emph{dollar delta} of the book on underlying $u$ is
\[
  \Delta^\$_u \;=\; \sum_{i : U(i) = u} n_i \, m \, \delta_i \, S_u,
\]
i.e.\ dollars gained per 1\% move in $S_u$.  Dollar delta is
additive across positions on the same underlying; aggregating
across underlyings is meaningful only if they are assumed to move
together (e.g.\ a beta-weighted aggregate).

\emph{Share-equivalent delta} $\sum_i n_i m \delta_i$ gives shares
to trade to hedge; \emph{dollar delta} gives loss per \% move.
Desk convention picks one.

\begin{keyfactbox}[Portfolio risk representation]
A book's risk is stored as a collection of \emph{bucketed Greeks},
not a single number per Greek:
\begin{itemize}
  \item \textbf{Delta} is bucketed by \emph{underlying} (one number
    per stock / index / currency pair).
  \item \textbf{Gamma} is bucketed by underlying and often by
    \emph{strike band} around spot.
  \item \textbf{Vega} is bucketed by underlying, \emph{tenor}
    (maturity band: 1M, 3M, 6M, 1Y, 2Y, 5Y\ldots) and often
    \emph{strike band} (25-delta put, ATM, 25-delta call), forming a
    full vega \emph{surface}.
  \item \textbf{Rho} is bucketed by currency and by \emph{tenor} on
    the rates curve (key-rate durations, KRDs).
\end{itemize}
Aggregation without bucketing discards hedging information.  A
``delta-neutral'' book can be long \$50M AAPL and short \$50M MSFT:
a pairs trade, not a hedged position.
\end{keyfactbox}

\subsection{Key-rate durations (KRDs)}

For rate-sensitive books (swaps, bonds, swaptions), a single
duration number is useless: you need to know \emph{which part of
the curve} you are exposed to.  The answer is the \emph{key-rate
duration} vector from \cref{ch:fixed_income_foundations}.

Let the zero curve be parameterised by pillar tenors $\{T_1, \ldots,
T_K\} = \{1\mathrm{M}, 3\mathrm{M}, 6\mathrm{M}, 1\mathrm{Y},
2\mathrm{Y}, 5\mathrm{Y}, 10\mathrm{Y}, 30\mathrm{Y}\}$.  The
key-rate duration at pillar $k$ is
\[
  \mathrm{KRD}_k
  \;\defeq\; -\frac{1}{V} \frac{\partial V}{\partial r_k},
\]
with other pillars held fixed.  Equivalently, bump $r_k$ by 1\,bp,
reprice, divide by $V \cdot 10^{-4}$.  The vector
$(\mathrm{KRD}_1, \ldots, \mathrm{KRD}_K)$ is the book's exposure
along the curve.  A 5Y-receive, 10Y-pay swap has near-zero total
duration but $\mathrm{KRD}_{5\mathrm{Y}} > 0$ and
$\mathrm{KRD}_{10\mathrm{Y}} < 0$: a pure curve-shape bet
(steepener).

KRDs satisfy a sum rule: on a parallel shift, $\sum_k
\mathrm{KRD}_k \cdot \Delta r_k = \mathrm{DV01}$ in appropriate
units.  They are additive across trades in the same currency.

\subsection{Cross-gamma}

On a book with multiple underlyings, a single gamma per underlying
misses the \emph{joint} convexity.  For a portfolio $V(S_1, S_2,
\ldots, S_n)$ the full second-order Taylor expansion is
\[
  \Delta V
  \;\approx\;
  \sum_i \Delta_i \, \Delta S_i
  \;+\; \tfrac{1}{2} \sum_i \sum_j \Gamma_{ij} \, \Delta S_i \Delta S_j,
\]
where $\Gamma_{ij} = \partial^2 V / (\partial S_i \partial S_j)$.
Diagonals $\Gamma_{ii}$ are the usual single-underlying gammas;
off-diagonals $\Gamma_{ij}$, $i \ne j$, are \emph{cross-gammas}:
the rate at which $\Delta_i$ changes as $S_j$ moves.

Cross-gamma is nonzero when the portfolio contains instruments
whose delta on one underlying depends on another's price.  Examples:
\begin{itemize}
  \item \emph{Index option} decomposed into constituents: a
    50-delta SPX call has delta $\approx 0.50$ against SPX, but
    across the 500 constituents each delta depends on the others
    via the index weight, producing large cross-gammas.
  \item \emph{Spread option} with payoff $\max(S_1 - S_2 - K, 0)$:
    $\partial^2 V / (\partial S_1 \partial S_2) < 0$, so rallying
    $S_2$ reduces delta to $S_1$.
  \item \emph{Basket option} on $w_1 S_1 + w_2 S_2$: cross-gamma
    $\ne 0$ whenever $w_1 w_2 \ne 0$.
\end{itemize}
Single-name books can ignore cross-gamma; index market-making books
cannot; ignoring it produces P\&L Explain residuals of several
percent.

\subsection{Worked example: a 4-position mini-book}
\label{sec:mini_book}

Four European options on four underlyings; $r = 4\%$, $m = 100$.

\begin{center}
\small
\begin{tabular}{@{}llrrrrr@{}}
\toprule
Underlying & Type & Qty & $S$ & $K$ & $\sigma$ & Tenor \\
\midrule
AAPL & Call & $+10$ & 150  & 150  & 28\% & 30d  \\
MSFT & Put  & $+5$  & 400  & 390  & 24\% & 60d  \\
GOOG & Call & $-8$  & 140  & 145  & 30\% & 45d  \\
SPX  & Call & $+20$ & 5000 & 5000 & 15\% & 90d  \\
\bottomrule
\end{tabular}
\end{center}

Summing per-contract Greeks times quantity times multiplier:

\begin{center}
\small
\begin{tabular}{@{}lrrr@{}}
\toprule
Position & Share delta & Dollar delta (\$) & Vega per 1\,vol\,pt (\$) \\
\midrule
AAPL $+10$ C 150  & $+532$  & $+79{,}847$    & $+171$ \\
MSFT $+5$  P 390  & $-177$  & $-70{,}661$    & $+301$ \\
GOOG $-8$  C 145  & $-326$  & $-45{,}654$    & $-153$ \\
SPX  $+20$ C 5000 & $+1{,}135$ & $+5{,}673{,}609$ & $+19{,}527$ \\
\midrule
\textbf{Total}    & $+1{,}164$ & $+5{,}637{,}141$ & (see below) \\
\bottomrule
\end{tabular}
\end{center}

The share-delta total is misleading: 97\% of the \$1{,}164 delta is
SPX at \$5000; AAPL/MSFT/GOOG are rounding error in dollar
terms.  Dollar delta makes this clear: almost \$5.7M is SPX.

A single vega total is useless: 30-day AAPL vol and 90-day SPX
vol respond to different moves.  The desk representation is a
\emph{vega bucket table} (rows = underlying, columns = tenor):

\begin{center}
\small
\begin{tabular}{@{}lrrrr@{}}
\toprule
Underlying & 30d & 45d & 60d & 90d \\
\midrule
AAPL & $+171.0$ & --- & --- & --- \\
MSFT & --- & --- & $+301.4$ & --- \\
GOOG & --- & $-152.7$ & --- & --- \\
SPX  & --- & --- & --- & $+19{,}527.0$ \\
\bottomrule
\end{tabular}
\end{center}

Each cell is dollars per 1 vol point move on that
$(underlying, tenor)$ bucket; a real desk sees a heatmap across
many more buckets.  Trader's question: ``if 30-day single-name
vols rally 2 points while index vol stays flat, what do I make?''
Answer: $2 \times (171 + 301.4 - 152.7) = +\$639$, dominated by
MSFT.  The SPX bucket is an order of magnitude larger than all
single-name vega combined: the book's vol risk lives there.

\begin{intuitionbox}
Aggregating Greeks without bucketing is like aggregating a
multi-currency P\&L without FX: the single number hides the
risk.  A ``delta-neutral'' book can be long one name and short
another; a ``vega-flat'' book can be long front-month vol and short
back-month, an enormous \emph{term-structure} bet.  The bucketed
view is the hedgeable view.
\end{intuitionbox}


%% ================================================================
\section{P\&L Explain: the Strat's daily proof-of-work}
\label{sec:pnl_explain}
%% ================================================================

Every morning on every derivatives desk, a Strat takes yesterday's
closing book plus yesterday's and today's market data, and
\emph{explains} the overnight change in mark-to-market (MTM) value
using the book's Greeks.  This report is \emph{P\&L Explain}
(or \emph{P\&L Attribution}), the Strat's daily proof-of-work.
If it balances, pricing library and risk system agree, the trader
starts the day, and Middle Office signs off on official P\&L.  If
not, the Strat's job is to find out why, fast.

\subsection{The Explain formula}

The overnight MTM change is a second-order Taylor expansion in the
market state variables:

\begin{keyfactbox}[P\&L Explain formula]
\begin{align*}
  \Delta V
  \;\approx\;\;
  & \underbrace{\Delta \cdot \Delta S}_{\text{delta P\&L}}
  \;+\; \underbrace{\tfrac{1}{2}\Gamma (\Delta S)^2}_{\text{gamma P\&L}}
  \;+\; \underbrace{\nu \, \Delta \sigma}_{\text{vega P\&L}} \\
  & \;+\; \underbrace{\Theta \, \Delta t}_{\text{theta P\&L}}
  \;+\; \underbrace{\rho \, \Delta r}_{\text{rho P\&L}}
  \;+\; \underbrace{\varepsilon}_{\text{residual}}.
\end{align*}
Healthy book: $|\varepsilon| < 1\%$ of $|\Delta V|$.  A residual
above a few percent flags a missing Greek, stale bump size, model
misalignment, or a late-booked trade.  Above 10\% means \emph{stop
trading until fixed}.
\end{keyfactbox}

Terms use \emph{start-of-day} Greeks (not end-of-day): the
Explain is a predictor of what P\&L \emph{should have been} given
the observed market moves.

For multi-underlying books, the delta, gamma, and cross-gamma terms
expand:
\[
  \text{2nd-order spot P\&L}
  \;=\; \sum_i \Delta_i \, \Delta S_i
  \;+\; \tfrac{1}{2} \sum_{i,j} \Gamma_{ij} \, \Delta S_i \Delta S_j.
\]
The vega term expands over the \emph{vol surface}: $\sum_{u,k,j}
\nu_{u,k,j} \, \Delta \sigma_{u,k,j}$ where $(u,k,j)$ indexes
(underlying, tenor, strike bucket).  The rho term expands over KRD
pillars: $\sum_k \rho_k \, \Delta r_k$.

\subsection{Decomposition: the canonical five bucket Explain}

The typical Explain groups the expansion into five buckets:

\begin{enumerate}
  \item \textbf{Delta / Spot P\&L.}  First-order spot move.
    Dominates for directional positions.
  \item \textbf{Gamma / Convexity P\&L.}  Second-order spot move.
    Positive for long-gamma books (market making), negative for
    short-gamma books (selling options).
  \item \textbf{Vega / Vol P\&L.}  Change in implied vols.
    Bucketed over the vol surface.
  \item \textbf{Theta / Time decay.}  Passage of time.  Usually
    negative for long-option books, positive for short-option.
  \item \textbf{Rho / Rate P\&L.}  Change in zero rates.  Usually
    small for options books; dominant for rates desks.
\end{enumerate}

Residual = MTM change minus the sum of these five.

\subsection{Worked example: ATM call, 1 day, $S{:}100\!\to\!101$,
$\sigma{:}20\%\!\to\!21\%$}
\label{sec:explain_worked}

A single ATM call on a non-dividend stock, start of day:
\[
  S_0 = 100, \quad K = 100, \quad r = 2\%, \quad \sigma_0 = 20\%,
  \quad T_0 = 0.25 \text{ (3 months)}.
\]
Black--Scholes price and Greeks (\cref{ch:greeks_hedging}), to 6
d.p.:
\[
  C_0 = 4.232160, \quad
  \Delta = 0.539828, \quad
  \Gamma = 0.039695,
\]
\[
  \nu = 19.847627 \text{ per 1.00 change in } \sigma, \quad
  \Theta = -8.934063 \text{ per year}, \quad
  \rho = 12.437656.
\]
Equivalently, $\nu = 0.198476$ per 1 vol point, $\Theta = -0.024477$
per calendar day.

By close of day 1:
\[
  S_1 = 101, \quad \sigma_1 = 21\%, \quad r_1 = r = 2\%,
  \quad T_1 = T_0 - 1/365 = 0.247260.
\]

\paragraph{Step 1: reprice.}
\[
  C_1 \;=\; 4.963545, \qquad
  \Delta V \;=\; C_1 - C_0 \;=\; 0.731386.
\]
Actual MTM change: \$0.731386 per share, \$73.14 per contract.

\paragraph{Step 2: contributions.}  Moves:
$\Delta S = +1$, $\Delta \sigma = +0.01$, $\Delta t = +1/365$,
$\Delta r = 0$.  Plug in:
\begin{align*}
  \text{delta P\&L} &= \Delta \cdot \Delta S
     = 0.539828 \times 1 = +0.539828, \\
  \text{gamma P\&L} &= \tfrac{1}{2}\Gamma (\Delta S)^2
     = 0.5 \times 0.039695 \times 1 = +0.019848, \\
  \text{vega P\&L} &= \nu \cdot \Delta \sigma
     = 19.847627 \times 0.01 = +0.198476, \\
  \text{theta P\&L} &= \Theta \cdot \Delta t
     = -8.934063 \times (1/365) = -0.024477, \\
  \text{rho P\&L} &= \rho \cdot \Delta r = 0.
\end{align*}
Sum: $0.539828 + 0.019848 + 0.198476 - 0.024477 + 0 = +0.733675$.

\paragraph{Step 3: residual.}
\[
  \varepsilon \;=\; \Delta V - \text{sum}
  \;=\; 0.731386 - 0.733675 \;=\; -0.002289.
\]
The residual is $-\$0.00229$ per share ($-0.31\%$ of $\Delta V$),
well inside the 1\% tolerance.  Explain balances.

\paragraph{Step 4: the residual.}  Sum of unmodelled
third-order Taylor terms:
\begin{itemize}
  \item \emph{Speed} $\partial \Gamma / \partial S$: higher-order
    convexity.
  \item \emph{Vanna} $\partial \Delta / \partial \sigma =
    \partial \nu / \partial S$: spot--vol cross-term.  With both
    $\Delta S$ and $\Delta \sigma$ nonzero, this dominates.
  \item \emph{Volga} $\partial^2 V / \partial \sigma^2$:
    second-order vol convexity.
  \item \emph{Charm} $\partial \Delta / \partial t$: how delta
    ages overnight.
\end{itemize}
On a market-making book with large vanna and volga (deep OTM
options), these terms contribute tens of basis points each and the
5-Greek Explain breaks down.  Fix: add \emph{second-order Greeks}
(vanna, volga, charm, speed), but that raises position cost from
five numbers to fifteen.  Trade-off: vanilla books stop at five;
exotic books carry the full second-order expansion.

Summary table:

\begin{center}
\small
\begin{tabular}{@{}lr@{}}
\toprule
Contribution & P\&L (\$ per share) \\
\midrule
Delta  & $+0.539828$ \\
Gamma  & $+0.019848$ \\
Vega   & $+0.198476$ \\
Theta  & $-0.024477$ \\
Rho    & $\phantom{-}0.000000$ \\
\midrule
Sum of Greeks     & $+0.733675$ \\
Actual $\Delta V$ & $+0.731386$ \\
\midrule
Residual          & $-0.002289$ \\
Residual / $|\Delta V|$ & $-0.31\%$ \\
\bottomrule
\end{tabular}
\end{center}

\begin{intuitionbox}
Clean P\&L Explain is the Strat's daily proof-of-work.  Sum of
Greeks matches MTM within 1\%: risk system and pricing library are
in sync, screen Greeks are trustworthy.  Residual above 1\%: a
missing vol bucket, stale bump size, late-booked trade, or pricing
change not propagated to risk.  Healthy books have residuals near
zero; unhealthy books have residuals drifting systematically,
usually a missing Greek (vega on a new bucket, cross-gamma on a new
basket).  A persistent 2\% residual eventually becomes a \$50M
mystery loss.
\end{intuitionbox}

\subsection{Reference implementation}

A minimal single-option P\&L Explain, matching the worked example
above:

\begin{lstlisting}[caption={Single-option P\&L Explain.}]
from math import log, sqrt, exp, pi, erf

def ncdf(x): return 0.5*(1 + erf(x/sqrt(2)))
def npdf(x): return exp(-x*x/2) / sqrt(2*pi)

def bs_call_greeks(S, K, r, sig, T):
    d1 = (log(S/K) + (r + 0.5*sig*sig)*T) / (sig*sqrt(T))
    d2 = d1 - sig*sqrt(T)
    price = S*ncdf(d1) - K*exp(-r*T)*ncdf(d2)
    delta = ncdf(d1)
    gamma = npdf(d1) / (S*sig*sqrt(T))
    vega  = S*npdf(d1)*sqrt(T)              # per 1.00 in sigma
    theta = (-S*npdf(d1)*sig/(2*sqrt(T))    # per year
             - r*K*exp(-r*T)*ncdf(d2))
    rho   = K*T*exp(-r*T)*ncdf(d2)          # per 1.00 in r
    return price, delta, gamma, vega, theta, rho

def pnl_explain(S0, K, r, sig0, T0, dS, dsig, dr, dt_days):
    C0, d, g, v, th, rh = bs_call_greeks(S0, K, r, sig0, T0)
    C1, *_            = bs_call_greeks(S0+dS, K, r+dr, sig0+dsig,
                                       T0 - dt_days/365)
    dt = dt_days / 365
    contribs = {'delta': d*dS,   'gamma': 0.5*g*dS*dS,
                'vega':  v*dsig, 'theta': th*dt, 'rho': rh*dr}
    total    = sum(contribs.values())
    residual = (C1 - C0) - total
    return C1 - C0, contribs, residual
\end{lstlisting}


%% ================================================================
\section{xVA: the alphabet soup of valuation adjustments}
\label{sec:xva}
%% ================================================================

Black--Scholes assumes the counterparty always pays, the dealer
always pays, the hedge is fundable at the risk-free rate, and no
capital is tied up against the position.  None of these is true in
practice.  The \emph{valuation adjustments} (``xVA'') are the price
of reality:

\begin{center}
\small
\begin{tabular}{@{}lll@{}}
\toprule
Adjustment & What it prices & Sign \\
\midrule
CVA & Counterparty default on positive exposures & $-$ (reduces value) \\
DVA & Own default on negative exposures         & $+$ (increases value) \\
FVA & Funding spread on uncollateralised positions & $-$ (typically) \\
KVA & Regulatory capital cost over life of trade & $-$ \\
MVA & Initial-margin funding (cleared trades)   & $-$ \\
\bottomrule
\end{tabular}
\end{center}

The risk-adjusted price is
\[
  V_{\mathrm{risky}}
  \;=\; V_{\mathrm{BS}}
  \;-\; \mathrm{CVA}
  \;+\; \mathrm{DVA}
  \;-\; \mathrm{FVA}
  \;-\; \mathrm{KVA}
  \;-\; \mathrm{MVA}.
\]
On a long-dated swap with a weak counterparty, xVA can run to 5\%
of notional: larger than the bid--ask spread, and the difference
between winning the trade and losing money.

\subsection{CVA: credit valuation adjustment}

CVA is the expected loss from counterparty default, the
derivative analogue of the bond formula $\mathrm{EL} =
(1-R)\,\mathrm{PD}$ from \cref{ch:credit_derivatives}, with
\emph{exposure} replaced by the stochastic MTM and integrated over
the trade's life.

\begin{keyfactbox}[CVA formula]
Let $V(t)$ be the dealer's MTM at time $t$.  On counterparty
default at time $\tau \le T$ the dealer's exposure is $V(\tau)^+$:
you lose only if in the money when they default.  Unilateral
CVA:
\[
  \mathrm{CVA}
  \;=\; (1-R) \int_0^T D(t)\,
       \E^\Qm \big[V(t)^+ \,\big|\, \tau = t\big]\,
       dPD(t),
\]
where $R$ is recovery, $D(t)$ the risk-free discount factor,
$\E[V(t)^+]$ the \emph{expected positive exposure} (EPE) at time
$t$, and $dPD(t) = -dS(t) = h(t) S(t)\,dt$ the marginal risk-neutral
default probability (\cref{ch:credit_derivatives}).
\end{keyfactbox}

Discrete form on grid $0 = t_0 < \cdots < t_N = T$:
\[
  \mathrm{CVA}
  \;\approx\; (1-R) \sum_{k=1}^{N}
  D(t_k) \cdot \mathrm{EPE}(t_k) \cdot
  \big[S(t_{k-1}) - S(t_k)\big],
\]
where the bracket is the marginal default probability in
$(t_{k-1}, t_k]$ stripped from CDS quotes.

\paragraph{Worked example: CVA on a 5-year interest-rate swap.}
Dealer is fixed receiver on a 5-year at-par swap; counterparty's
5-year CDS trades at 150\,bp ($R = 40\%$).  Flat hazard $h =
s/(1-R) = 0.025$ gives survival $S(t) = e^{-0.025 t}$; risk-free
rate $r = 3\%$.  EPE on an at-par vanilla swap has the hump
$\mathrm{EPE}(t) \approx N \sqrt{t(T-t)/T}$, peaking near $T/2$;
take $N = 100$.  Semi-annual steps:

\begin{center}
\small
\begin{tabular}{@{}rrrrr@{}}
\toprule
$t$ (yr) & $\mathrm{EPE}(t)$ & $D(t)$ & $\Delta S_k$ & contribution \\
\midrule
0.5 &  67.08 & 0.9851 & 0.01242 & 0.4925 \\
1.0 &  89.44 & 0.9704 & 0.01227 & 0.6389 \\
1.5 & 102.47 & 0.9560 & 0.01212 & 0.7121 \\
2.0 & 109.54 & 0.9418 & 0.01196 & 0.7406 \\
2.5 & 111.80 & 0.9277 & 0.01182 & 0.7354 \\
3.0 & 109.54 & 0.9139 & 0.01167 & 0.7010 \\
3.5 & 102.47 & 0.9003 & 0.01152 & 0.6379 \\
4.0 &  89.44 & 0.8869 & 0.01138 & 0.5417 \\
4.5 &  67.08 & 0.8737 & 0.01124 & 0.3953 \\
5.0 &   0.00 & 0.8607 & 0.01110 & 0.0000 \\
\midrule
\multicolumn{4}{r}{\textbf{Sum} $= $} & \textbf{5.60} \\
\bottomrule
\end{tabular}
\end{center}

Each contribution is $(1-R) \cdot D(t_k) \cdot \mathrm{EPE}(t_k)
\cdot \Delta S_k$.  CVA is \$5.60 per \$100 of peak EPE, or 5.6\%
of peak exposure.  On a \$1bn notional swap (with a calibrated EPE
replacing the simplified hump), CVA runs into single-digit
millions, a charge the salesperson must quote on top of the
BS-implied swap rate.

\subsection{DVA: debit valuation adjustment}

DVA is the mirror of CVA: the benefit the \emph{dealer} gets from
their own default probability.  If I owe you \$100M and default, I
pay only $R \cdot 100$M, so ex ante my expected payout is less
than \$100M, the liability is worth less than \$100M to me, and my
balance sheet gains when the market perceives my credit
deteriorating.

DVA uses the dealer's own survival curve and \emph{negative}
exposure:
\[
  \mathrm{DVA}
  \;=\; (1-R_{\mathrm{own}}) \int_0^T D(t)\,
       \E^\Qm \big[(-V(t))^+\big]\,
       dPD_{\mathrm{own}}(t).
\]

\begin{pitfallbox}[DVA is accounting-real but not cash-real]
DVA is the most controversial xVA.  Under IFRS 13 and US GAAP
(post-2007), banks must recognise DVA in fair-value accounting,
so when a bank's own CDS widens in a crisis, it books a \emph{gain}
on its derivative liabilities.  Hence the Q3 2011 headlines where
Morgan Stanley, Goldman Sachs, and Citi reported multi-billion DVA
gains as their own spreads blew out, gains they could never
monetise unless they defaulted.  Many banks now strip DVA from
``core'' earnings.  Economic logic clean; political optics
terrible.
\end{pitfallbox}

\subsection{FVA: funding valuation adjustment}

When a dealer's uncollateralised derivative becomes an asset
(positive MTM), the dealer must borrow cash to hedge at their
unsecured rate, which exceeds the risk-free rate.  The difference
is the \emph{funding spread}; FVA is its discounted integral over
the trade's life.

Symmetrically, if the position becomes a liability, cash received
can be invested at the unsecured rate, a funding benefit (FBA).
FVA is often quoted net of FBA.

Approximate formula:
\[
  \mathrm{FVA}
  \;\approx\; \int_0^T D(t) \cdot s_{\mathrm{fund}}(t)
              \cdot \E^\Qm\big[V(t)^+\big]\, dt,
\]
where $s_{\mathrm{fund}}(t)$ is the dealer's funding spread over
OIS.  Same EPE integral as CVA, different weight (funding spread
vs.\ hazard $\times$ LGD).

FVA applies only to \emph{uncollateralised} trades.  Fully
collateralised trades (CSA-perfect daily cash margining at OIS)
have zero FVA because the collateral funds the hedge, which is
why mandatory clearing post-Dodd--Frank (\cref{ch:regulatory})
dramatically reduced FVA on vanilla flow.  (A \emph{CSA}, Credit
Support Annex, is the ISDA document specifying daily collateral
exchange; \emph{OIS}, the overnight-index-swap rate, is the
near-risk-free rate paid on that posted cash.)

\subsection{KVA: capital valuation adjustment}

Under Basel III, banks hold regulatory capital against derivative
exposures: CVA-risk capital (capital against CVA itself moving),
default-risk capital (RWA $\times$ 8\%), plus trading-book
market-risk capital.  Shareholders demand (say) 10\% return on
equity.  Over a long-dated trade this return on capital is
non-trivial and is passed to the client as KVA.

\[
  \mathrm{KVA}
  \;\approx\; \int_0^T D(t) \cdot
              \mathrm{ROE} \cdot \mathrm{Cap}(t)\, dt,
\]
where $\mathrm{Cap}(t)$ is regulatory capital consumed at time $t$
and ROE is target return on equity.  On a 30-year uncollateralised
swap with a low-rated counterparty, KVA can be the largest xVA.

\subsection{MVA: margin valuation adjustment}

MVA is the cost of funding \emph{initial margin} posted to a CCP
(central counterparty, the clearing house that interposes itself
between trade counterparties) or bilaterally under uncleared
margin rules.  IM is typically funded at the dealer's unsecured
rate but earns only OIS inside the CCP, producing a spread cost.
Less commonly quoted than CVA/DVA/FVA, increasingly important
post-UMR (Uncleared Margin Rules, phased 2016--2022, requiring
two-way IM on bilateral derivatives).

\begin{keyfactbox}[xVA alphabet soup]
\begin{itemize}
  \item \textbf{CVA}: counterparty default.  Negative to dealer;
    uses counterparty hazard, EPE.
  \item \textbf{DVA}: own default benefit.  Reported in IFRS/GAAP,
    often stripped from adjusted earnings.
  \item \textbf{FVA}: funding cost on uncollateralised positions;
    zero for trades collateralised at OIS.
  \item \textbf{KVA}: regulatory capital $\times$ ROE over life;
    dominant on long-dated, low-rated, uncollateralised trades.
  \item \textbf{MVA}: initial-margin funding spread.
\end{itemize}
On a weak-credit 30y uncollateralised swap: CVA $\sim$ 100\,bp,
FVA $\sim$ 50\,bp, KVA $\sim$ 200\,bp.  Total xVA can exceed the
BS bid--ask by 5$\times$.
\end{keyfactbox}

\begin{pitfallbox}[xVA is not naively additive across the book]
Novices compute per-trade CVA/FVA and sum.  Wrong in two ways.

\textbf{Netting.}  Under an ISDA Master with CSA, counterparty
exposure is \emph{netted} across all trades.  A \$100M ITM swap
and a \$90M OTM swap with the same counterparty net to \$10M, not
\$100M; per-trade CVA over-estimates tenfold.

\textbf{Collateral.}  Variation margin reduces CVA to the
margin-period-of-risk tail ($\sim$10 days).  Fully collateralised
in cash at OIS: CVA and FVA near zero; only MVA and KVA remain.

xVA is computed at \emph{netting-set} level; attributing to
individual trades (for incremental pricing) is a separate
calculation called \emph{incremental xVA}.
\end{pitfallbox}


%% ================================================================
\section{Risk limits and VaR}
\label{sec:risk_limits}
%% ================================================================

Bucketed Greeks and xVA tell a desk where its risk lives;
\emph{limits} constrain how much is allowed.  A typical desk
operates under a hierarchy:

\begin{itemize}
  \item \textbf{Greek caps.}  Hard caps on dollar delta, share
    delta, and vega per underlying and vega bucket, e.g.\ ``AAPL
    30-day vega capped at \$50k per vol point.''
  \item \textbf{Notional caps.}  Gross and net notional by product
    and counterparty.
  \item \textbf{Scenario P\&L caps.}  Loss limit under a stress,
    e.g.\ ``5\% SPX drop with vols up 20\% must not exceed
    \$10M.''
  \item \textbf{VaR and Expected Shortfall.}  Statistical tail
    measures of loss.
\end{itemize}

\subsection{VaR and Expected Shortfall}

\emph{Value-at-Risk} (VaR) at confidence $\alpha$ and horizon $h$
is the loss not exceeded with probability $\alpha$ over $h$:
\[
  \mathrm{VaR}_\alpha(h)
  \;=\; \inf \{ \ell \ge 0 : \Prob(L_h \le \ell) \ge \alpha \},
\]
with $L_h$ the loss.  Typical: $\alpha = 99\%$, $h = 1$ day.

\emph{Expected Shortfall} (ES, also CVaR) is expected loss given a
tail event:
\[
  \mathrm{ES}_\alpha(h)
  \;=\; \E[L_h \mid L_h \ge \mathrm{VaR}_\alpha(h)].
\]
ES is always at least as large as VaR at the same confidence and
is a \emph{coherent} risk measure (VaR is not: it can violate
subadditivity, penalising diversification).

The Basel Fundamental Review of the Trading Book (FRTB) replaced
VaR with 97.5\% ES as the regulatory market-risk measure from 2023
(\cref{ch:regulatory}).

\subsection{Why a single VaR number is not enough}

VaR/ES is a \emph{portfolio-level aggregate}: it collapses the
risk distribution to one number.  Too coarse for a trader running
the book in real time: it does not say \emph{which} Greek drives
risk, or whether a breach is vega or delta.  Desks run VaR
\emph{plus} a dashboard of bucketed Greeks.  VaR answers the CRO's
question (``how bad could tomorrow be?''); bucketed Greeks answer
the trader's (``what do I hedge now?'').

\begin{intuitionbox}
Risk limits are the desk's exoskeleton.  Greeks tell you
\emph{where} risk is; limits tell you \emph{whether} you are
allowed to have it.  VaR/ES is the regulator's view, bucketed
Greeks the trader's.  Both exist because neither alone suffices:
ES without Greeks is un-hedgeable; Greeks without ES miss tail
correlations.  In production both are computed nightly and
monitored intraday.
\end{intuitionbox}


%% ================================================================
\section{Key facts to memorise}
%% ================================================================

\begin{itemize}
  \item Portfolio Greeks are \emph{bucketed} by underlying, tenor,
    strike band, and curve pillar, not aggregated into scalars.
  \item \textbf{Dollar delta} $= n m \delta S$ gives \$ per \% move;
    \textbf{share delta} $= n m \delta$ gives shares-to-hedge.
  \item KRDs decompose rate risk by pillar; parallel shifts recover
    DV01 as $\sum_k \mathrm{KRD}_k \Delta r_k$.
  \item Cross-gamma $\Gamma_{ij} = \partial^2 V /(\partial S_i
    \partial S_j)$ is essential for multi-underlying books (index
    options, spreads, baskets).
  \item \textbf{Explain formula}: $\Delta V \approx \Delta \cdot
    \Delta S + \tfrac12 \Gamma (\Delta S)^2 + \nu\, \Delta\sigma +
    \Theta\, \Delta t + \rho\, \Delta r + \varepsilon$; healthy
    residual $|\varepsilon| < 1\%$ of $|\Delta V|$.
  \item Residual is third-order Taylor: speed, vanna, volga, charm.
    When vanna/volga matter, add second-order Greeks.
  \item \textbf{CVA} $= (1-R) \int_0^T D(t) \mathrm{EPE}(t)\, dPD(t)$,
    with hazard stripped from CDS (\cref{ch:credit_derivatives}).
  \item \textbf{DVA} is the symmetric own-default adjustment,
    controversial because it books a gain when your own credit
    deteriorates.
  \item \textbf{FVA} is funding spread on uncollateralised exposures;
    zero for perfectly OIS-collateralised trades.
  \item \textbf{KVA} = ROE $\times$ regulatory capital over life;
    dominant on long-dated, low-rated, uncollateralised trades.
  \item xVA is computed at \emph{netting-set} level; netting and
    collateral can reduce CVA by an order of magnitude.
  \item Risk-adjusted price: $V_{\mathrm{risky}} = V_{\mathrm{BS}}
    - \mathrm{CVA} + \mathrm{DVA} - \mathrm{FVA} - \mathrm{KVA}
    - \mathrm{MVA}$.
  \item VaR is not coherent (violates subadditivity); ES is.  FRTB
    uses 97.5\% ES.
  \item A desk runs bucketed Greeks, stress scenarios, and VaR/ES
    in parallel; each answers a different question.
\end{itemize}


\begin{prosperitybox}[title={Prosperity example: per-round P\&L Explain is trivial}]
Prosperity has no vol surface, rate curve, or counterparty, so the
multi-Greek Explain collapses to one line: ``end-of-round cash
$+$ end-of-round inventory at mid $-$ start cash $=$ round P\&L.''
The discipline is still universal: after each round, split P\&L
into \emph{executed-trade} (realised) and \emph{mark-to-market
inventory} (unrealised), attribute each line to a strategy (market
making, stat arb, options scalping), and reconcile to the
leaderboard.  If attribution does not match the posted number, the
local backtester is wrong, and an unexplained \$100/round gap
becomes a \$2000-notional bug on submission day.  The same
instinct that chases a 50\,bp Explain residual on a sell-side desk
saves a Prosperity team from submitting a broken algo.
\end{prosperitybox}

\begin{gsbox}
On a Goldman Sachs equity-derivatives or rates-desk Strat seat,
the 7am--9am workflow has been fixed for two decades.
\textbf{07:00.}  Overnight batch is done; pull previous-close
book, previous-close data, today's open data.  \textbf{07:15.}
Run P\&L Explain, the morning's most important job; report
goes to Market Risk, Product Control, and trader by 08:30.
\textbf{07:30.}  If Explain balances ($< 1\%$ residual), move to
xVA checks: any counterparty's CDS moved enough to shift CVA?  Any
new trade reprice the netting set?  \textbf{08:00.}  Investigate
stale-trade or missing-Greek exceptions from the overnight batch.
\textbf{08:30.}  Send consolidated Explain and xVA report.
\textbf{09:00.}  Market open; pivot to trader support: pricing
requests, scenario runs, hedge simulations.  Same cadence on every
derivatives desk in every major bank.  An Explain that does not
balance on Friday is why Strats work weekends.
\end{gsbox}

\begin{historybox}[title={Historical note: from accounting afterthought to trillion-dollar business}]
Pre-2008, CVA was a quarterly back-office adjustment with minimal
trader involvement.  Lehman's September 2008 default (and
Bear Stearns' March rescue) made counterparty default a
first-order risk overnight.  Basel III (2010--2011) introduced
\emph{CVA-risk capital} (explicit capital against CVA itself
moving), turning CVA into a first-class trading activity.  Banks
set up dedicated \emph{CVA desks} that hedge counterparty credit
by trading CDS, pricing CVA into every long-dated trade.

FVA was more contentious.  Hull and White (2012)
\citep{hullwhite2012fva} argued FVA should not enter fair value
because it is a financing cost, not a counterparty cost.  The
industry disagreed: JPMorgan took a \$1.5bn FVA charge in 2014 and
every major dealer followed.  The accounting debate is technically
open; market practice is settled.
\end{historybox}
